% Appendix B — Numerical Analysis Primer
\chapter{Numerical Analysis Primer}
\label{app:numerical-primer}

This appendix reviews the numerical analysis foundations assumed by the main
text, providing self-contained derivations for readers whose background is
primarily in physics rather than computational mathematics.

\section{Error Estimation and Convergence Order}
\label{sec:num-error}

We define truncation error, discretisation error, and the concept of
convergence order, and show how Richardson extrapolation can be used to
estimate the error from a sequence of grid resolutions.

\section{The CFL Condition}
\label{sec:num-cfl}

We derive the Courant-Friedrichs-Lewy stability condition for explicit
time-stepping schemes, relating the time step to the grid spacing and the
fastest propagation speed in the system.

\section{Lax Equivalence Theorem}
\label{sec:num-lax}

We state the Lax-Richtmyer equivalence theorem --- that for a consistent
finite-difference scheme, stability is necessary and sufficient for
convergence --- and discuss its implications for numerical relativity.

\section{Floating-Point Arithmetic}
\label{sec:num-floating-point}

We review the IEEE 754 floating-point representation, define machine epsilon
and unit roundoff, and discuss catastrophic cancellation and its mitigation
through compensated summation and the double-single technique used in
Chapter~\ref{ch:ray-tracing}.

\paragraph{Key References.}
Error analysis follows \cite{LeVeque2007}; the CFL condition follows
\cite{Courant1928}; floating-point arithmetic follows \cite{Higham2002};
the Lax equivalence theorem follows \cite{Lax1956,Richtmyer1967}.
